\section{Resolution Road Map and Evidence Upgrading}

\subsection{A staged program rather than a single limit}

The question of pointlike behavior is best treated as a staged program. The first stage is synthetic closure. It asks whether the selected channel, binning, response model, and nuisance ledger can recover a known injected pointlike, contact, form-factor, or excited-state pattern. This is a test of the inference apparatus, not a stand-in for a collision measurement. It is especially valuable because a fitted contact coefficient can otherwise be driven by a poorly modeled tail or a hidden response correlation.

The second stage is channel qualification. A channel enters the common comparison only after its kinematic definition, fiducial selection, theory calculator, response model, and data status are registered. A public binned measurement without a response interface can be included as a qualitative evidence source but not as a quantitative likelihood. A public likelihood can be used only under its documented validity conditions. A future collider or electron-ion facility remains a projection until it produces released data. This staged admission protects the program from assigning equal weight to incomparable inputs.

The third stage is cross-channel synthesis. The combined analysis should not ask whether every channel has the same apparent deviation. Instead, it asks whether a particular physics model predicts a consistent coefficient, form-factor scale, or excited-state benchmark after the channels retain their own PDFs, backgrounds, and response effects. A disagreement may rule out the selected common model, expose an interface error, or point to a more complicated alternative. It does not become evidence for compositeness by itself.

\begin{table*}[!t]
\caption{Resolution-ladder evidence-upgrading rules.}
\label{tab:roadmap}
\centering
\footnotesize
\begin{tabularx}{\textwidth}{p{0.15\textwidth}p{0.23\textwidth}Y Y}
\toprule
\textbf{Stage} & \textbf{Admissible input} & \textbf{Primary question} & \textbf{Go or no-go rule}\\
\midrule
Closure & Synthetic records generated from stated hypotheses and nuisance models & Can the inference recover the injected class and separate it from adversarial nuisance shifts? & No-go for quantitative interpretation if classes are non-identifiable.\\
Qualification & Released binning, selections, response information, and theory provenance & Is the channel definition sufficient to construct a reproducible likelihood? & No-go for combined fit if status or response provenance is missing.\\
Single-channel fit & Pointlike and alternative predictions under a declared cutoff policy & Does any alternative improve held-out prediction without a nuisance explanation? & Report only a channel-conditioned preference or constraint.\\
Cross-channel synthesis & Qualified channels connected by an explicit model relation & Is a shared alternative compatible with all relevant channel likelihoods? & No-go for a common interpretation if coefficient or shape relations conflict.\\
Independent confirmation & A separate observable or facility with materially different systematics & Does the preferred pattern survive a new resolution and uncertainty environment? & Upgrade only after independent confirmation and repeated validity tests.\\
\bottomrule
\end{tabularx}
\end{table*}

\subsection{How each outcome changes the scientific answer}

The pointlike-SM branch is not an empty outcome. If it survives the qualified likelihood comparison, it supports a sharper statement: within the registered momentum-transfer interval, fiducial definition, uncertainty model, and alternative class, no additional fermion substructure is required. This statement can be compared across runs because the resolution ledger records the terms under which it is true. It remains stronger than a casual assertion that nothing was found, yet it does not outgrow its evidence.

The contact-EFT branch has a different interpretation. A stable contact coefficient can identify a departure from the pointlike baseline in an effective description, but it cannot by itself identify a constituent scale or a unique ultraviolet completion. The analysis must disclose whether the preferred coefficient is robust to changes in cutoff, dimension-eight terms, and flavor assumptions. High-energy tail studies make this issue explicit: the utility of the tail is matched by sensitivity to the domain in which the EFT description is reliable \cite{allwicher2023}.

The form-factor branch asks a still different question. A radius-like parameter can summarize a smooth momentum-transfer dependence in a valid low-$Q$ expansion. Its apparent value must be tested against response functions and PDF variations that can also tilt a spectrum. An excited-fermion branch adds resonance, acceptance, and reconstruction questions. Existing excited-lepton searches illustrate why a benchmark limit must retain its mass and coupling assumptions \cite{cms_tau2025}. RLMC keeps those hypotheses separate because their failure modes are separate.

\subsection{Strategic value of complementary facilities}

Energy, luminosity, flavor reach, and experimental systematics are not interchangeable resources. A future hadron collider can extend high-$Q$ reach; a lepton collider can emphasize cleaner initial conditions and precision; an electron-ion program can constrain partonic structure with a different kinematic organization. Their complementarity is valuable only when it is expressed through a transfer condition and status label, not through a decorative collection of sensitivity plots. Planning studies for future colliders and electron-ion measurements provide motivation for these distinct roles \cite{abada2019,arrington2021}.

The resulting scientific answer is dynamic but rigorous. The smallest currently tested fields are those whose pointlike-SM descriptions remain adequate after the best available resolution tests and stated alternatives. A future result may lower a radius bound, exclude an excited-state region, restrict an EFT coefficient, or reveal a coherent deviation. Each possibility refines the answer without requiring the proposal to claim that it has already occurred.
